\section{The original local wave packet and its global normalized product}
\label{sec:bessel-wave-packets-v49}

\status{This chapter proves the whole original \texttt{CoefWeightedLoc}
and \texttt{WlocExists}.  The latter has the original
\texttt{BesselNVGap} antecedent.  It also proves the original
\texttt{WFromLoc} and, using the completed setting blocks on their
original supplied \texttt{Forms0} scope, constructs
\texttt{ConcreteSettingExists} conditional on that
nonvanishing input.  No proof of \texttt{BesselNVGap}, Bessel
multiplicity one, or the remaining Siegel--Weil comparisons is asserted.
The construction keeps every supplied Haar and quotient measure.
Its global vector is chosen before the later measured block.}

All Hilbert inner products in this chapter are conjugate linear in the
first argument.  Put \(n=2r+k+1\), \(R=r+1\), and
\[
 N=N(k,r)=4n^2\sqrt{2R}.
 \tag{BWP.1}
\]
The source's \(H_{k,r}\), after reversing the order of the final
isotropic basis, is the matrix
\[
 H_{k,r}=\begin{pmatrix}0&0&I_r\\0&J_V&0\\I_r&0&0\end{pmatrix},
 \qquad J_V=\operatorname{diag}(I_k,-1).
 \tag{BWP.2}
\]
This reordering is unitary, so it does not change the source operator
norm \(o(g)=\|g\|_{\mathrm{op}}\).  Write \(G=U(H_{k,r})\),
\(N_0=\texttt{Nset}(k,r)\), \(H_z=\texttt{UWset}(k,r,z)\),
and \(S=H_zN_0\).  The symbol \(N_0\) denotes the group and
\(N\) the exponent.  The original \(z\) belongs to the middle
space and satisfies \(-z^*J_Vz>0\).

The analytic foundations used below are Haar existence and uniqueness,
the quotient Haar theorem with its Borel-section formulation, Tonelli
and Bochner integration, and finite-dimensional compact representation
theory.  We also retain the same Harish--Chandra admissibility theorem
used in RFD1: every irreducible unitary representation of a real
reductive group has finite-dimensional compact isotypic spaces.
Borel's \emph{Automorphic forms on reductive groups}, printed
Section 10, paragraph labelled 8.5, states this precise input.
The reproducing function itself is proved below, not taken as a further
analytic theorem.  Source predicates and numerical assumptions are
recorded in the accompanying contracts and source excerpts.

\subsection{Explicit unipotent coordinates and compact reproduction}

% BEGIN BWP1
\paragraph{BWP1. A global chart on the actual unipotent subgroup.}
\label{claim:BWP1}
Let \(A\) run through upper unitriangular complex \(r\)-by-\(r\)
matrices, \(B\) through complex \(r\)-by-\((k+1)\) matrices,
and \(Z\) through skew-Hermitian \(r\)-by-\(r\) matrices.
Then
\[
 \Phi(A,B,Z)=
 \begin{pmatrix}
 A&B&(Z-\tfrac12BJ_V^{-1}B^*)A^{-*}\\
 0&I&-J_V^{-1}B^*A^{-*}\\
 0&0&A^{-*}
 \end{pmatrix}
 \tag{BWP.3}
\]
is a polynomial homeomorphism of a real vector space of dimension
\(d=2r^2+2rk+r\) onto the literal \(N_0\).

\begin{proof}
The source fixes the successive isotropic flag quotients pointwise,
and fixes the quotient of its orthogonal complement by that flag.
Its defining column equations therefore give the block shape
\([A,B,C;0,I,E;0,0,F]\), with \(A\) upper unitriangular.
The equation \(g^*H_{k,r}g=H_{k,r}\) first gives
\(F=A^{-*}\) and \(E=-J_V^{-1}B^*A^{-*}\).
The remaining equation is
\(CA^*+AC^*=-BJ_V^{-1}B^*\).
Thus \(Z=CA^*+\tfrac12BJ_V^{-1}B^*\) is skew-Hermitian,
and solving for \(C\) gives (BWP.3).  Conversely those equations
verify unitarity and every defining column condition.
The inverse reads \(A,B\) from their blocks and uses the displayed
quadratic formula for \(Z\).  Since the inverse of a unitriangular
matrix is a finite polynomial, both maps are polynomial.
The three real dimensions are \(r(r-1),2r(k+1),r^2\).
Their sum is the stated \(d\).  When \(r=0\) the chart is the
one-point zero-dimensional chart, and the same proof applies.
\end{proof}
% END BWP1

% BEGIN BWP2
\paragraph{BWP2. Haar measure and reciprocal-norm integrability in that chart.}
\label{claim:BWP2}
Lebesgue measure in (BWP.3) is Haar measure on \(N_0\).
Every measure satisfying the source's weak \texttt{IsHaarOn}
on \(N_0\) is a finite positive multiple of it, and
\[
 \int_{N_0}o(m)^{-N}\,d\mu_{N_0}(m)<\infty.
 \tag{BWP.4}
\]
This proves the original \texttt{NIntDecay} at all \(k,r\).

\begin{proof}
For a fixed left multiplier \((A_0,B_0,Z_0)\), direct block
multiplication gives
\[
 \begin{split}
 A'&=A_0A,\qquad B'=A_0B+B_0,\\
 Z'&=A_0ZA_0^*+Z_0+
 \tfrac12(A_0BJ_V^{-1}B_0^*-B_0J_V^{-1}B^*A_0^*).
 \end{split}
 \tag{BWP.5}
\]
The real Jacobian is block triangular.  The first diagonal block
has determinant one by ordering upper entries row by row; the second
has determinant \(|\det A_0|^{2(k+1)}=1\); congruence by \(A_0\)
on skew-Hermitian matrices has determinant of absolute value
\(|\det A_0|^{2r}=1\).  Thus (BWP.5) preserves Lebesgue measure.
AMD8 promotes any given weak invariant measure to Haar, so uniqueness
gives a finite positive scalar multiple of this measure.

The matrix in (BWP.2) is a self-adjoint unitary involution.  Hence
\(g^{-1}=H_{k,r}g^*H_{k,r}\), \(o(g^{-1})=o(g)\), and
\(o(g)\ge1\) on \(G\).  Every entry of a matrix has absolute
value at most its operator norm.  The inverse-chart formula in BWP1
therefore gives, for any fixed Euclidean coordinate norm,
\[
 1+\|(A-I,B,Z)\|\le C_{k,r}\,o(\Phi(A,B,Z))^2.
 \tag{BWP.6}
\]
For \(r>0\), \(d< n^2\), whereas \(N\ge4n^2\); in
particular \(N/2>d\).  Raising (BWP.6) to exponent \(-N/2\)
majorizes the integrand in (BWP.4) by a constant times
\((1+\|t\|)^{-N/2}\), which is integrable in real dimension
\(d\) by radial integration.  For \(r=0\), the Haar measure
of the one-point group is finite and the integral is immediate.
\end{proof}
% END BWP2

% BEGIN BWP3
\paragraph{BWP3. The actual Bessel subgroup and its unitary character.}
\label{claim:BWP3}
The literal groups satisfy \(S=N_0\rtimes H_z=H_zN_0\), with
unique factorization in either order.  The group \(H_z\) is compact,
\(S\) is closed and unimodular, and
\(\psi=\texttt{psiGen}(k,r,z,c)\) is a unitary character of
\(N_0\) invariant under conjugation by \(H_z\).  If \(\pi\)
has the original \texttt{IsHCTypeW}, then
\[
 \rho(hm)=\psi(m)\pi(h)\quad(h\in H_z,\ m\in N_0)
 \tag{BWP.7}
\]
is a continuous unitary representation of \(S\).

\begin{proof}
An element of \(H_z\) fixes the two isotropic blocks, acts on the
middle space by an element of \(U(k,1)\), and fixes its negative
vector \(z\).  Its restriction to \(z^\perp\) identifies this
group with \(U(k)\).  Reading the induced middle block of an element
of \(S\) gives its unique \(H_z\) factor, and the remaining
factor has exactly the columns in BWP1.  These readings are continuous;
equivalently the flag and middle-action equations show closedness.
The same argument is the local form of AMD4--AMD6.  The unipotent
group is unimodular, and the action modulus of the compact group
\(H_z\) on its Haar measure is a continuous homomorphism into
\(\mathbb R_{>0}\), hence one.  The semidirect product is unimodular.

For indices below the last one, the source function \(\ell_a(m)\)
is the \((a,a+1)\) entry of \(A\); for the last index it is
the last coordinate of \(Bz\).  Formula (BWP.5) shows that all
these functions add under multiplication: the final row of \(A_0\)
is the final coordinate row.  Conjugation by \(H_z\) fixes these
entries because it fixes the isotropic basis and \(z\).
Consequently \(\exp(i\operatorname{Re}\sum_a c_a\ell_a)\)
is the stated invariant unitary character, including the empty sum
when \(r=0\).  The defining frame of \texttt{IsHCTypeW}
transports a unitary finite-dimensional \(U(k)\) representation
to the whole \(H_z\).  This proves every assertion about (BWP.7).
\end{proof}
% END BWP3

% BEGIN BWP4
\paragraph{BWP4. A compact reproducer in every prescribed neighborhood.}
\label{claim:BWP4}
Let \(\sigma\) be any original \texttt{IsDS} representation,
and let \(w\) belong to its lowest compact-type subspace.  For
every neighborhood \(U\) of the identity there is
\(\alpha\in C_c(G,\mathbb C)\), supported in \(U\), such that
\[
 \int_G\alpha(a)\sigma(a)w\,d\mu_G(a)=w.
 \tag{BWP.8}
\]
This proves the source's \texttt{HCReprodDS}, and its weaker
compact measurable versions, without any gap assumption.

\begin{proof}
The source \texttt{IsDS} includes a complete Hilbert space and a
continuous irreducible unitary representation.  The given vector is
compact finite.  Harish--Chandra admissibility therefore makes the sum
\(E\) of its full compact isotypic spaces finite dimensional.  This
uses the named inherited foundation stated above, not a consequence
of square integrability alone.  If \(w=0\), take \(\alpha=0\).
Otherwise put \(m=\dim E\).  Compactness of the maximal compact
group \(K\) gives a conjugation-invariant neighborhood \(V\) of
one, containing one, with \(V^m\subset U\).
Take a continuous nonnegative approximate-identity function supported
in a sufficiently small subset of \(V\), of integral one, and
average its conjugates over \(K\).  Conjugation preserves Haar
measure.  The resulting \(f\in C_c(G)\) is conjugation invariant,
has integral one and support in \(V\).  Strong continuity, uniformly
on the unit sphere of the finite-dimensional \(E\), makes
\(T=\sigma(f)|_E\) arbitrarily close to the identity.  It preserves
\(E\), because \(\sigma(f)\) commutes with \(K\), and it is
invertible for a sufficiently small choice.

Write its characteristic polynomial as
\(p(t)=a_0+a_1t+\cdots+a_mt^m\); here \(a_0\ne0\).
Cayley--Hamilton gives
\(I_E=-a_0^{-1}\sum_{j=1}^m a_jT^j\).
Thus \(\alpha=-a_0^{-1}\sum_{j=1}^m a_j f^{*j}\) has
the required action.  Each convolution is continuous and compactly
supported, with support in \(V^j\subset V^m\subset U\).
Bochner integrability follows from compact support and unitarity.
The same argument works for any supplied Haar normalization.
\end{proof}
% END BWP4

\subsection{The exact coefficient exponent}

% BEGIN BWP5
\paragraph{BWP5. The original cubic gap dominates every required analytic bound.}
\label{claim:BWP5}
Suppose \(\lambda=\texttt{lam48}(k,r,v,b)\) is an original
Harish--Chandra parameter and \(\texttt{GapBig}(k,r,b)\) holds.
Choose its witness \(t\), and put
\[
 M=n+(2r+2)t,\qquad h=2M^3+t+3.
 \tag{BWP.9}
\]
Every opposite-sign parameter difference has absolute value at least
\(h\).  If \(q_0=\min(p,q)>0\), then
\[
 h>npq+1,\quad h>16nq_0+8\sqrt{2q_0},\quad
 h>16nq_0+8N\sqrt{2q_0}.
 \tag{BWP.10}
\]

\begin{proof}
The first \(k-1\) signs are positive; the remaining signs alternate.
Their parameters are \(-b_j\).  Each consecutive difference of
the \(b_j\)'s is strictly larger than \(h\), so each distinct
tail difference is larger than \(h\).  A negative tail index has
odd tail index and therefore lies after the zeroth tail entry.
Strict decrease of \(\lambda\) places every initial parameter
above that zeroth entry; its distance to a negative tail entry is
again greater than \(h\).  These are all opposite-sign pairs.
Counting negative tail indices gives \(q_0\le q\le R\).

Since \(t\ge n+5\), we have \(M\ge3n\) and \(M\ge10R\).
Consequently \(npq\le M^3\) and
\(96n^2R\le2(3n)^2(10R)\le2M^3<h\).
Using \(N=4n^2\sqrt{2R}\), we obtain
\[
 16nq_0+8N\sqrt{2q_0}
 \le16nR+64n^2R\le80n^2R<h.
\]
Also \(\sqrt{2q_0}\le2R\), so
\(16nq_0+8\sqrt{2q_0}\le16nR+16R\le32nR<h\).
Finally \(2M^3+t+3>M^3+1\).  These prove (BWP.10),
including all small values permitted by the original predicates.
\end{proof}
% END BWP5

% BEGIN BWP6
\paragraph{BWP6. The whole original weighted coefficient statement.}
\label{claim:BWP6}
For every original \(k,r,v,b,\sigma,w\) in
\texttt{CoefWeightedLoc}, one has, with its prescribed real Haar measure,
\[
 \int_G |\langle w,\sigma(g)w\rangle|\,o(g)^N\,dg<\infty.
 \tag{BWP.11}
\]
No additional regularity or model-identification premise is required.

\begin{proof}
If \(q_0=0\), the real unitary group is compact and the integrand
is continuous, so the assertion follows.  Otherwise BWP5 supplies
the full DC.2 bounds.  DC13 therefore identifies the given literal
\texttt{IsDS} with the constructed representation, including its
prescribed lowest compact type.  The operator estimate DC6 holds
for every unit vector in that type, hence for the given \(w\).
Writing \(g=k\exp X\), \(a=\|X\|\), it bounds the
coefficient by
\(\exp[-(h/8)\min(a^2,a/\sqrt{2q_0})]\).
The operator norm satisfies \(o(g)\le e^a\), and DC6 bounds the
polar Jacobian by \(e^{n\sqrt{2q_0}a}\), times a polynomial
radial shell factor.  For large \(a\) the total exponential rate is
\[
 -h/(8\sqrt{2q_0})+N+n\sqrt{2q_0}<0
\]
by the final inequality in (BWP.10).  Radial integration is finite;
near zero continuity suffices.  Haar rescaling changes the answer
by a finite positive constant.  This proves the literal universal
\texttt{CoefWeightedLoc}, not merely an estimate for one chosen model.
\end{proof}
% END BWP6

% BEGIN BWP7
\paragraph{BWP7. Transport to the given hyperbolic coordinates and exact kernel.}
\label{claim:BWP7}
Retain the original \texttt{lam48}/\texttt{eps48} parameters and
the \texttt{IsHCParam}/\texttt{GapBig} hypotheses of BWP5--BWP6.
For the original \texttt{IsDSHpk} and \texttt{IsLowUnitH} data
at these parameters,
every supplied weak Haar measure \(\mu_G\) is admissible in
(BWP.8) and (BWP.11).  Put
\[
 d_\sigma=\frac{\|w\|^4}
 {\int_G|\langle w,\sigma(g)w\rangle|^2\,d\mu_G(g)},
 \qquad e(g)=d_\sigma\overline{\langle w,\sigma(g)w\rangle}.
 \tag{BWP.12}
\]
Then \(d_\sigma>0\), this is precisely the source
\texttt{fdegH}/\texttt{eFun}, and
\[
 \int_G e(g)\sigma(g)w\,d\mu_G(g)=w,
 \qquad \int_G|e(g)|o(g)^N\,d\mu_G(g)<\infty.
 \tag{BWP.13}
\]

\begin{proof}
AMD8 first promotes \(\mu_G\) to genuine Haar measure on \(G\).
The Cayley matrix in each defining predicate is Euclidean unitary.
Conjugation by it preserves operator norms exactly and sends \(G\)
to the specified diagonal unitary group.  If the two predicates supply
different Cayley matrices, their ratio is Euclidean unitary and
preserves the same diagonal signature, hence belongs to its maximal
compact group.  Inner conjugation by that element preserves every
compact isotypic space and its minimum-norm type.  Thus the given
\(w\) and \(\sigma\) can use the same Cayley coordinates.
The transported measure is a finite positive multiple of the prescribed
Haar measure.  BWP4 and BWP6 now transport back with the exact original
\(\mu_G\).

DC12, and DC13 for the given representation, give Schur orthogonality
for this measure.  With \(\|w\|=1\), its formal degree is exactly
the quotient in (BWP.12).  The weighted integrability follows from
BWP6.  The Bochner integral in (BWP.13) exists by unitarity and
unweighted integrability.  Pairing it with any vector \(v\), Schur
orthogonality evaluates the result to \(\langle v,w\rangle\).
The integral is therefore \(w\).  This checks the conjugation and
formal-degree normalization in the literal source kernel.
\end{proof}
% END BWP7

\subsection{The packet for the original quotient measure}

% BEGIN BWP8
\paragraph{BWP8. The original quotient height and support.}
\label{claim:BWP8}
Let \(\mu_G,\mu_S,\nu\) be exactly the measures in
\texttt{WlocExists}.  Then \(\nu\) is supported on \(G\).
On \(G\), the source height
\[
 \eta(x)=\inf_{s\in S}o(sx)=\texttt{nuH}(k,r,z,x)
 \tag{BWP.14}
\]
is Borel measurable, satisfies \(1\le\eta(x)\le o(x)\),
and is left \(S\)-invariant.  The original quotient identity can
therefore be used with this exact height, without choosing a new \(\nu\).

\begin{proof}
Both weak measures are genuine Haar by AMD8.  Apply the given
quotient identity to the indicator of the complement of \(G\).
Multiplication by \(s\in S\subset G\) preserves that complement.
The inner integral equals \(\mu_S(S)>0\) on it and zero on
\(G\), so its \(\nu\)-measure is zero.  The norm inequalities
follow from BWP2 and \(1\in S\).  The infimum of continuous
functions is upper semicontinuous: its strict sublevel sets are the
unions of their strict sublevel sets.  Finiteness and the lower bound
thus give Borel measurability.  Replacing \(s\) by \(ss_0\)
proves \(\eta(s_0x)=\eta(x)\).  All statements concern the
source ambient measure, even when it is given by some other choice
of representatives.
\end{proof}
% END BWP8

% BEGIN BWP9
\paragraph{BWP9. The exact product scalar for the two Bessel measures.}
\label{claim:BWP9}
Choose any Haar measures \(\mu_H,\mu_N\) on the actual
\(H_z,N_0\), viewed on the ambient matrix space.  Their product
under \((h,m)\mapsto hm\) is \(c_S\mu_S\) for a finite
positive real number \(c_S\).  Thus, for every nonnegative
measurable \(f\),
\[
 \int\!\int_{H_z\times N_0} f(hmx)\,d\mu_H\,d\mu_N\,d\nu(x)
 =c_S\int_G f(g)\,d\mu_G(g).
 \tag{BWP.15}
\]
The same conclusions hold for any weak Haar measures quantified by
the original nonvanishing premise.

\begin{proof}
BWP3 gives the semidirect product and its trivial action modulus.
Consequently the product is Haar, as in AMD7; uniqueness relative to
the supplied \(\mu_S\) gives exactly the scalar stated.  Substitute
that identity into the supplied quotient formula.  AMD8 applies to
weak Haar measures and makes \(\mu_H\) finite because \(H_z\)
is compact, and \(\mu_N\) finite on compact sets.  No rescaling of
the given \(\mu_G,\mu_S,\nu\) has been made.
\end{proof}
% END BWP9

For the next claims fix the complete original real numerical setting,
the given \(\sigma,\pi,w\), and the measures above.  Let
\(x_0\in G\), \(\xi\in V_\pi\), and set
\[
 F(g)=\langle\sigma(x_0)w,\sigma(g)w\rangle,\qquad
 A(x)=\int_{H_z\times N_0}
 F(hmx)\overline{\psi(m)}\,\pi(h)^{-1}\xi\,d\mu_H(h)d\mu_N(m).
 \tag{BWP.16}
\]
These are the source's \texttt{Af} and its exact Bessel integrand.
On matrices outside \(G\), choose the value zero for the constructed
functions; all measures in the assertions are supported on \(G\).

% BEGIN BWP10
\paragraph{BWP10. Absolute convergence, continuity, and a uniform polynomial bound.}
\label{claim:BWP10}
The integral (BWP.16) converges absolutely for every \(x\in G\),
is continuous there, and satisfies
\[
 \int_{H_z\times N_0}\|F(hmx)\overline{\psi(m)}\pi(h)^{-1}\xi\|
 \,d\mu_Hd\mu_N\le C\,o(x)^N
 \tag{BWP.17}
\]
for a finite constant depending on the fixed data and \(x_0,\xi\).

\begin{proof}
Unitarity gives \(F(g)=\langle w,\sigma(x_0^{-1}g)w\rangle\).
By BWP7 and left translation,
\(L=\int_G|F(y)|o(y)^N\,d\mu_G(y)<\infty\).
Choose the compact reproducer \(\alpha\) from BWP4 for \(w\),
and let \(K\) be its compact support.  Applying its vector identity
to the coefficient gives
\(F(y)=\int_K\alpha(a)F(ya)\,d\mu_G(a)\).
For \(y=hmx\), set \(u=hmxa\).  Since
\(m=h^{-1}ua^{-1}x^{-1}\), BWP2 gives
\[
 o(m)\le o(h^{-1})o(u)o(a^{-1})o(x).
\]
The first and third factors are bounded uniformly on \(H_z\times K\).
Bound \(|\alpha|\) by its supremum and change variable from
\(a\) to \(u\) using left Haar invariance.  This proves
\[
 |F(hmx)|\le C_0L\,o(x)^N o(m)^{-N}.
 \tag{BWP.18}
\]
The remaining factors in the integrand have norm \(\|\xi\|\).
BWP2 and finiteness of \(\mu_H\) prove (BWP.17).
On every compact set of \(x\)'s, (BWP.18) is a single integrable
majorant.  Strong continuity of the representations and dominated
convergence prove continuity of \(A\).  This argument justifies
absolute convergence before every later change of variables.
\end{proof}
% END BWP10

% BEGIN BWP11
\paragraph{BWP11. The literal equivariance of the packet.}
\label{claim:BWP11}
For all \(h\in H_z,m\in N_0,x\in G\),
\[
 A(hmx)=\psi(m)\pi(h)A(x).
 \tag{BWP.19}
\]
In particular \(\|A\|\) is left \(S\)-invariant.

\begin{proof}
By BWP9 the integral is the same fixed scalar times
\(\int_S F(sx)\rho(s)^{-1}\xi\,d\mu_S(s)\).
For \(s_0\in S\), replace \(s\) by \(s's_0^{-1}\).
Right invariance, supplied by BWP3, and
\(\rho(s's_0^{-1})^{-1}=\rho(s_0)\rho(s')^{-1}\)
give \(A(s_0x)=\rho(s_0)A(x)\).
Absolute convergence is BWP10, and unitarity in BWP3 gives the norm
assertion.  Formula (BWP.19) is precisely \texttt{IsSPsiEquiv}.
\end{proof}
% END BWP11

% BEGIN BWP12
\paragraph{BWP12. Weighted integrability for the supplied quotient measure.}
\label{claim:BWP12}
The packet satisfies
\[
 \int_G\|A(x)\|\eta(x)^N\,d\nu(x)<\infty,
 \qquad \int_G\|A(x)\|\,d\nu(x)<\infty.
 \tag{BWP.20}
\]

\begin{proof}
Bound the norm of (BWP.16) by the integral of its norm.  For every
\(h,m\), BWP8 gives \(\eta(x)\le o(hmx)\).
Tonelli, followed by (BWP.15), therefore bounds the first integral by
\[
 c_S\|\xi\|\int_G |F(g)|o(g)^N\,d\mu_G(g),
\]
which is finite by the translated estimate in BWP10.  This uses only
nonnegative integration, so no unproved integrability is used to
exchange the integrations.  The second assertion follows from
\(\eta\ge1\).  Continuity and BWP8 supply the measurability needed.
\end{proof}
% END BWP12

% BEGIN BWP13
\paragraph{BWP13. Square integrability from the quotient height.}
\label{claim:BWP13}
The source squared norm satisfies
\(\texttt{l2SG}(\nu,A)=\int\|A\|^2d\nu<\infty\).

\begin{proof}
For every \(s\in S\), BWP10 and BWP11 give
\(\|A(x)\|=\|A(sx)\|\le C o(sx)^N\).
Taking the infimum, and using continuity and monotonicity of the
positive power, gives \(\|A(x)\|\le C\eta(x)^N\).
Multiply by \(\|A(x)\|\), integrate, and use BWP12.
The conclusion does not require a uniform bound on \(A\) over
the quotient; the original exponent supplies exactly the needed bound.
\end{proof}
% END BWP13

% BEGIN BWP14
\paragraph{BWP14. What the original nonvanishing input supplies.}
\label{claim:BWP14}
Assume the original \texttt{BesselNVGap}.  There are
\(x_0,x\in G\) and \(\xi\in V_\pi\) for which the packet
(BWP.16) satisfies
\[
 A(x)\ne0,\qquad 0<\texttt{l2SG}(\nu,A)<\infty.
 \tag{BWP.21}
\]
In explicit terms, the retained premise says, for every original
numerical setting, representations, lowest unit vector and pair of
weak Haar measures, that some \(x_0,x\in G\) make
\[
 \int_{H_z\times N_0}
 \langle\sigma(x_0)w,\sigma(hm)\sigma(x)w\rangle
 \overline{\psi(m)}\,\pi(h)^{-1}\,d\mu_Hd\mu_N\ne0
 \tag{BWP.30}
\]
as an endomorphism of the given \(V_\pi\).  This endomorphism
nonvanishing is not proved here; the asserted packet positivity is
its consequence.

\begin{proof}
The original real setting has \(1\le k\) and \(k\) even,
hence \(2\le k\).  All its remaining parameter, interlacing,
half-integral difference, gap, and nonzero-character conditions are
exactly those quantified by \texttt{BesselNVGap}.  Apply it to the
given representations, given lowest unit vector, and the measures
chosen in BWP9.  It gives a nonzero endomorphism integral for some
\(x_0,x\).  A nonzero endomorphism has a vector \(\xi\) with
nonzero image.  Evaluation on that vector is (BWP.16); evaluation
commutes with its absolutely convergent finite-dimensional integral.

Finiteness of the squared norm is BWP13.  To prove positivity for the
actual \(\nu\), choose a nonnegative nonzero compactly supported
continuous function \(f\) on an open set where \(\|A\|>\delta>0\).
Such a set exists by BWP10, and \(\int_G f\,d\mu_G>0\)
by full support of Haar measure.  If the squared norm were zero,
then \(A=0\) almost everywhere for \(\nu\).  Equivariance
implies \(f(sy)=0\) for all \(s\) whenever \(A(y)=0\).
The supplied quotient identity would give \(\int_G f=0\), a
contradiction.  This proves positivity without replacing the quotient
measure or presupposing full support of a chosen section.
\end{proof}
% END BWP14

% BEGIN BWP15
\paragraph{BWP15. A justified convolution transfer for every base point.}
\label{claim:BWP15}
Suppose a measurable scalar function \(\beta\) satisfies
\[
 \int_G|\beta(g)|o(g)^N\,d\mu_G(g)<\infty,
 \qquad \int_G\beta(g)\sigma(g)w\,d\mu_G(g)=w.
 \tag{BWP.22}
\]
Then for every \(x\in G\) the following integral is absolutely
convergent and
\[
 \int_G\beta(g)A(xg)\,d\mu_G(g)=A(x).
 \tag{BWP.23}
\]

\begin{proof}
BWP10 and submultiplicativity give the bound on the integral of the
norm of the joint integrand
\[
 \int_G |\beta(g)|
 \int_{H_z\times N_0}|F(hmxg)|\,\|\xi\|\,d\mu_Hd\mu_Nd\mu_G
 \le C o(x)^N\int_G|\beta(g)|o(g)^N\,d\mu_G<\infty.
\]
All joint functions are measurable by continuity on the supported
groups and the assumed measurability of \(\beta\).
Bochner Fubini therefore applies.  For fixed \(h,m\), the inner
integral is obtained by applying the bounded linear map
\[
 v\longmapsto
 \langle\sigma(x_0)w,\sigma(hm)\sigma(x)v\rangle
 \overline{\psi(m)}\pi(h)^{-1}\xi
\]
to the vector integral in (BWP.22).  It gives the integrand of
\(A(x)\).  This proves (BWP.23) simultaneously at every
\(x\in G\), with the stated absolute integrability.
\end{proof}
% END BWP15

% BEGIN BWP16
\paragraph{BWP16. Reproduction by the literal formal-degree kernel.}
\label{claim:BWP16}
For every \(x\in G\), the source \texttt{eFun} convolution
with \(A\) is integrable and equals \(A(x)\).

\begin{proof}
Take \(\beta=e\) from (BWP.12).  Both hypotheses of BWP15
are exactly BWP7.  Thus its conclusion is
\(\texttt{rconv}(\mu_G,\texttt{eFun}(\mu_G,\sigma,w),A,x)=A(x)\),
with the full integrability clause.  In particular the same formal
degree is computed from the given \(\mu_G\), rather than borrowed
from a differently normalized group measure.
\end{proof}
% END BWP16

% BEGIN BWP17
\paragraph{BWP17. Compact reproduction of the packet.}
\label{claim:BWP17}
There is a bounded measurable \(\alpha\), supported in a compact
set \(\Omega\subset G\) containing one, such that for every
\(x\in G\) its convolution with \(A\) is integrable and
equals \(A(x)\).

\begin{proof}
Use BWP4 and BWP7 to choose the continuous compact reproducer of
\(w\) for the actual \(\mu_G\).  A continuous compactly supported
function is bounded and has finite integral after multiplication by
the continuous weight \(o(g)^N\).  Adjoin one to its compact support
to obtain \(\Omega\), and extend by zero on the ambient matrix
space.  Closedness of \(G\) gives measurability of that extension.
BWP15 proves both assertions about the convolution at every base point.
\end{proof}
% END BWP17

% BEGIN BWP18
\paragraph{BWP18. The entire original local packet producer.}
\label{claim:BWP18}
The whole original \texttt{WlocExists} holds.  Namely, under its
literal \texttt{BesselNVGap} and \texttt{CoefWeightedLoc}
antecedents, and for every supplied local numerical setting,
representation, lowest unit vector, and triple of measures, there
exists \(W\) having all seven original conclusions.

\begin{proof}
Choose \(x_0,\xi\) by BWP14 and put
\[
 L_2=\int_G\|A(x)\|^2\,d\nu(x),\qquad W=L_2^{-1/2}A.
 \tag{BWP.24}
\]
BWP13--BWP14 give \(0<L_2<\infty\), so the positive real
normalizing scalar is defined.  Continuity is BWP10; literal
\(S,\psi,\pi\) equivariance is BWP11; integrability and reproduction
by the exact \texttt{eFun} is BWP16.  BWP12 supplies both ordinary
and \(\texttt{nuH}^{N1H}\)-weighted \(L^1\) integrability.
The squared \(L^2\) integral is
\(L_2^{-1}L_2=1\), precisely the source \texttt{l2SG}, rather
than an unnormalized Hilbert norm.  BWP17 supplies the measurable,
bounded, compactly supported reproducer, with support enlarged to
contain one.  Scalar multiplication preserves every convolution and
equivariance equality.  These are all the original clauses, on the
same \(\mu_G,\mu_S,\nu\).  BWP6 in addition independently supplies
the coefficient antecedent, so nonvanishing is the only unproved
local input left by this construction.
\end{proof}
% END BWP18

\subsection{Assembly on the given global data}

% BEGIN BWP19
\paragraph{BWP19. The quotient product formula with its actual scalar.}
\label{claim:BWP19}
For a finite product \(G=\prod_vG_v\), \(S=\prod_vS_v\)
of the preceding unimodular groups, fix the original product group
measure and any supplied \(S\) Haar measure and quotient measure
\(\nu_Q\).  Choose local Haar measures \(\mu_{S,v}\) and
their normalized quotient measures \(\nu_v\).  There is one
finite positive scalar \(C\), independent of the test functions,
such that for every family of nonnegative measurable
\(S_v\)-invariant functions \(f_v\),
\[
 \int_G\prod_v f_v(x_v)\,d\nu_Q(x)
   = C\prod_v\int_{G_v} f_v(x_v)\,d\nu_v(x_v).
 \tag{BWP.25}
\]

\begin{proof}
Haar uniqueness gives \(\mu_S=a\prod_v\mu_{S,v}\),
\(0<a<\infty\).  The product of the local quotient measures
satisfies the quotient formula for \(\prod_v\mu_{S,v}\),
by Tonelli.  Thus the constant is \(C=a^{-1}\).
For clarity, the given ambient \(\nu_Q\) need not coincide
as a measure with the product of chosen section measures.
What is unique is its pushforward to \(S\backslash G\).
Here is the normalization argument.  The quotient Haar theorem
used in ASG17 gives, for this closed unimodular subgroup, a
nonnegative Borel cutoff \(b\) with
\(\int_S b(sx)\,d\mu_S(s)=1\).
It can also be constructed from a locally finite cover of the
second-countable quotient, compactly supported functions having
positive orbit averages on that cover, division by those averages,
and a partition of unity.  For any nonnegative invariant \(f\),
the supplied quotient formula applied to \(bf\) identifies
\(\int f\,d\nu_Q\) with \(\int_G bf\,d\mu_G\).
The product quotient formula gives
\(\int_G bf\,d\mu_G=a^{-1}\int f\,d(\prod_v\nu_v)\).
This proves equality on all such
functions, including infinite integrals.  Taking \(f=\prod_v f_v\)
and using Tonelli proves (BWP.25).
\end{proof}
% END BWP19

% BEGIN BWP20
\paragraph{BWP20. Every supplied global frame has the required local packet.}
\label{claim:BWP20}
Fix the original blocks \(A,L,P,R_f,B_u,\Pi,M_u\) and
\(\texttt{GapHyp}(D)\).  Suppose the local numerical gaps and
\texttt{BesselNVGap} hold.  For each original word \(u\) and
real place \(v\), BWP18 gives a packet in the actual supplied
space \(\Pi.V_{\mathrm{loc}}(u,v)\), for its actual compact
representation and lowest vector, with quotient measure \(\nu_v\).
It has the original local character, convolution, weighted integral,
and compact reproduction after transport to the given frame.

\begin{proof}
Use the given vectors \(e_i,f_i\), positive orthonormal frame
\(w_j\), and the vector \(z_u/\sqrt{T_u}\) to form the
matrix \(Q_v\).  Their original pairing identities say exactly
\(Q_v^*J_{\epsilon_v}Q_v=H_{k,r}\), so \(Q_v\) is
invertible and conjugation identifies the two unitary groups.
The same computation identifies the literal flag stabilizer,
unipotent subgroup and compact factor in BWP3.  The supplied compact
representation becomes exactly \(\texttt{IsHCTypeW}\) in these
coordinates: its matrix on the positive frame is the middle positive
block of the conjugated compact element.

There are two different coordinate changes here.  Choose a
Euclidean unitary Cayley matrix \(P_0\) with
\(P_0^*J_{\epsilon_v}P_0=H_{k,r}\), obtained by replacing each
hyperbolic pair by its normalized sum and difference, followed by
the required permutation of signs.  Define
\[
 C_G(g)=P_0^{-1}gP_0,\qquad C_S(s)=Q_v^{-1}sQ_v,
 \qquad T(x)=Q_v^{-1}xP_0.
 \tag{BWP.28}
\]
The representation used for the local packet is
\(\sigma_B(M)=\sigma_v(P_0MP_0^{-1})\), on the given Hilbert
space with the unchanged \(w_v\).  By this unitary Cayley matrix
it has the literal \texttt{IsDSHpk} and \texttt{IsLowUnitH}.
The generally nonunitary frame matrix \(Q_v\) is used for the
subgroup and the base point, not to change that lowest vector.
Direct multiplication gives the two identities
\[
 T(sx)=C_S(s)T(x),\qquad T(xg)=T(x)C_G(g).
 \tag{BWP.29}
\]

More explicitly, the signature is \(\texttt{eps48}(k,r)\);
the parameter is \((v_0,\ldots,v_{k-2},-b_0,\ldots,-b_{2r+1})\).
The compact highest weight has the two original interlacing fields.
They are precisely \texttt{InterlacingCondition}.  Their parities
give \texttt{HalfIntegralDifference}.  The vector becomes
\(\sqrt{T_u}\) times the final negative standard vector.
For the character take every \(c_a=4\pi/g_D\).
The frame identities give
\(\ell_a(Q_v^{-1}sQ_v)=\mathrm{rootCoord}_a(s)\).
Consequently its exponent is
\[
 i\operatorname{Re}\sum_a(4\pi/g_D)\ell_a
 =\frac{2\pi i}{g_D}\sum_a
       (\mathrm{rootCoord}_a+\overline{\mathrm{rootCoord}_a}),
\]
which is exactly the original \texttt{psiR}.  Since \(g_D>0\),
every \(c_a\ne0\).  This verifies the actual numerical hypotheses
of the local producer, including its character normalization.

Choose local subgroup and quotient measures with the prescribed
\(\mu_{G,v}\), as in ASG17.  Push the group measure by \(C_G\),
the subgroup measure by \(C_S\), and the quotient measure by \(T\).
They satisfy the exact quotient formula: (BWP.29) gives its right
side, while \(T=L_{Q_v^{-1}P_0}\circ C_G\) gives its left side
by Haar left invariance.  Thus BWP18 applies to exactly this triple.
If its packet is \(W_B\), set \(W_v(x)=W_B(T(x))\).
The first identity in (BWP.29) gives equivariance and the second gives
the convolution formula with the given \(w_v\).  The formal-degree
integral under \((C_G)_*\mu_{G,v}\) is literally the original
one by change of variables, so its scalar is unchanged.
Both quotient integrals, including their normalization, pull back
under \(T\).  The compact reproducer pulls back under \(C_G\),
which is a group homeomorphism preserving one.

Finally, put \(c_v=\max\{1,\|Q_v\|\|P_0^{-1}\|\}\).
Writing \(y=T(x)\) and \(s_B=C_S(s)\) gives
\(sx=Q_vs_ByP_0^{-1}\), and hence
\(\eta_v(x)\le c_v\eta_B(T(x))\).
This finite positive constant transfers the exact weighted estimate
with exponent \(N\).  Continuity follows from the displayed maps.
No compact representation, lowest vector, Hilbert space, or supplied
group measure has been replaced.
\end{proof}
% END BWP20

% BEGIN BWP21
\paragraph{BWP21. The normalized global vector on the supplied tensor data.}
\label{claim:BWP21}
For the data of BWP20 there is a \(W_b:\texttt{WBlk}\) with
every literal \texttt{IsWu} field, on the actual \(\nu_Q(u)\)
and the given tensor map, chosen independently of any later
\(M_b:\texttt{MBlk}\).

\begin{proof}
Fix a word and apply BWP19 to the supplied subgroup measure,
which is the product of its given \(N\) and \(H\) measures by
AMD7--AMD8.  Let \(C>0\) be its scalar in (BWP.25).  The field
has a real place; fix one \(v_0\).  Multiply the local packet at
\(v_0\) by \(C^{-1/2}\), leaving every other one unchanged.
Call these functions \(W_v\), and define
\[
 W_u(x)=\operatorname{tens}_u\bigl((W_v(x_v))_v\bigr).
 \tag{BWP.26}
\]
The supplied tensor-inner-product identity gives
\(\|W_u(x)\|=\prod_v\|W_v(x_v)\|\).
It also implies multilinearity of the given tensor map: pair an
additivity or scalar-linearity difference in one coordinate with
every pure tensor.  The inner-product product formula makes each
pairing zero, and the supplied spanning property makes the difference
zero.  The product norm identity then gives continuity of this
finite multilinear map.
The finite tensor product is complete and finite dimensional,
as in the original representation block.

Every \(s\) in the actual stabilizer factors place by place as
\(h_vm_v\).  The supplied global representation is trivial on
the unipotent factors and agrees with the tensor product of the given
compact representations.  Multiplying (BWP.19) over the places
therefore gives the literal global character-equivariance identity.
The formal-degree kernel is exactly the product of the local kernels.
Each local convolution is absolutely integrable by BWP18.  Tensor
multilinearity, its product norm identity, and finite-product Fubini
show that the global convolution is integrable and equals (BWP.26).
This checks the original kernel, including its formal degrees.

Continuity follows from continuity of each factor.  Formula (BWP.25)
applied to the product of their invariant norms gives finite \(L^1\).
Applied to their squared norms it gives
\(C\cdot C^{-1}\cdot\prod_v1=1\).
These are all five \texttt{IsWu} fields.  The constructed local
functions also give the exact \texttt{Wloc}, continuity, and
\texttt{Wu\_tens} fields of \texttt{WBlk}.
Make these choices for each original word.  The construction used
only the original preceding blocks and their measures; it did not
use a later adelic measured block.
\end{proof}
% END BWP21

% BEGIN BWP22
\paragraph{BWP22. Decay of that same global vector for every later measured block.}
\label{claim:BWP22}
The \(W_b\) chosen in BWP21 satisfies the entire original
\texttt{WuDecay} for every later \(M_b:\texttt{MBlk}\).
In particular no uniqueness theorem for Bessel models is needed
for this existential constructor.

\begin{proof}
Let \(\eta_v(x_v)=\inf_{s_v\in S_v}\|s_vx_v\|\), and
let \(\eta_\infty\) be the original \texttt{nuInf}.
Independent choices of the finitely many \(s_v\)'s give
\(\eta_\infty(x)\le\prod_v\eta_v(x_v)\).
All factors are at least one in the diagonal unitary coordinates.
The original exponent is
\(N_X=4n^2\sqrt{2\max_v q_v}\le N\), since
\(q_v\le r+1\) by the sign count in BWP5.
The norm comparison in BWP20 therefore implies finite
\(\int\|W_v\|\eta_v^{N_X}d\nu_v\), using the local
\(N\)-weighted estimate and the harmless finite factors
\(c_v^{N_X}\).  Formula (BWP.25), the tensor norm identity, and
the preceding infimum inequality now give
\[
 \int\|W_u(x)\|\eta_\infty(x)^{N_X}\,d\nu_Q(x)<\infty.
 \tag{BWP.27}
\]
This is the full original \texttt{WuWeighted} clause.

For each place take the compact reproducer supplied by BWP18 after
transport, and put \(\alpha_\infty((g_v))=\prod_v\alpha_v(g_v)\).
It is measurable, bounded, and supported in the compact product
\(\prod_v\Omega_v\), which contains one.  The local integrability
clauses and tensor norm identity justify Fubini; their reproducing
equalities give the global one at every \(x\).  This is exactly
\texttt{WuReprod}.  Neither clause uses \(M_b\): the original
\(\nu_Q\), \(W_u\), infinite group measure and exponents belong
to the already chosen blocks.  They consequently hold for every
later \(M_b\), in the required order \(\exists W_b\,\forall M_b\).
This proves decay for the constructed vector; it asserts nothing
about uniqueness or every arbitrary \texttt{IsWu} vector.
\end{proof}
% END BWP22

% BEGIN BWP23
\paragraph{BWP23. The original global gap supplies every local cubic gap.}
\label{claim:BWP23}
The full original \texttt{GapBigOfGapHyp}(D) holds, for every
original dimension identification and every real place.

\begin{proof}
Write \(n=D.R.r=2r+k+1\) and \(m=2r+2\).  The original
rank conditions give \(n\) odd and \(m\le n\).
The source \texttt{GapHyp} says \(n\ge1801\) and every
successive \(b\)-difference is at least \(16n^{19}\).
Set \(t=n+6\).  It is odd, is at least \(n+5\), and
\(n+t\ge1801\).  Since \(n\ge7\),
\(M=n+mt\le n+n(n+6)\le2n^2\).
It follows that
\[
 2M^3+t+3\le16n^6+n+9<16n^{19}.
\]
Thus each of the required strict cubic-gap inequalities follows from
the original weak lower bound.  The source \(b_{48}\) is the
same \(b\)-tuple reindexed by the supplied dimension equality;
transporting a finite index changes none of these inequalities.
This proves the entire original predicate, with no new gap assumption.
\end{proof}
% END BWP23

% BEGIN BWP24
\paragraph{BWP24. The exact remaining condition for the complete analytic setting.}
\label{claim:BWP24}
The original \texttt{WFromLoc}(D) holds for every \(D\).
Using the already reconstructed blocks for an original \(D\)
with its supplied \(\texttt{Forms0}(D,D.u_0)\), the retained
\texttt{BesselNVGap} implies
\texttt{ConcreteSettingExists}(D).  The original \texttt{Forms0}
scope of the lattice constructor is preserved.  The nonvanishing statement,
and the later spectral, orbital and arithmetic comparisons, remain
separate producers.

\begin{proof}
Under the original antecedents of \texttt{WFromLoc}, BWP20--BWP22
construct the required \(W_b\) for every supplied sequence of
blocks, with decay for every later \(M_b\).  BWP23 supplies the
local gaps from the global one.  These are exactly the quantifiers
and conclusion of \texttt{WExists}; the local packet input is
BWP18, and its coefficient input is BWP6.

For the full setting on that supplied \texttt{Forms0} scope,
choose the group instance, lattice, additive
character, reference frame, transport, bundle, compact representation
and quotient measures by ASG1--ASG17, in their original order.
Apply the preceding constructor to choose \(W_b\).  Only then
apply AMD12 to obtain \(M_b\) and one bridge with the literal
\texttt{MeasTie}.  BWP22 gives decay for precisely that measured
block.  Finally the infinite projector kernel is a finite product
of the kernels in (BWP.12).  BWP6 and \(o(g)\ge1\) give local
absolute integrability; product Fubini gives
\(\operatorname{Integrable}(X.\mathrm{ph},\mu_\infty)\).
Thus the source \texttt{ConcreteSettingExists} conclusion, including
one setting, one bridge, the exact ties, decay and kernel
integrability, follows.
No step proved the nonzero Bessel endomorphism required in BWP14.
The conclusion therefore retains \texttt{BesselNVGap} and the
stated original \texttt{Forms0} input to the lattice construction.
Nor does a completed analytic setting alone
identify the orbital measures or prove the original Siegel--Weil
and common-family counting statements.
\end{proof}
% END BWP24
